% ==========================================
%   الگوی خالی برای هر تمرین/پروژه‌ی جدید
%   این فایل را کپی کنید و به steps/stepXX.tex تغییر نام دهید
% ==========================================
\section[عنوان کوتاه برای فهرست مطالب]{\faProjectDiagram\ \ عنوان کامل پروژه/تمرین}

\begin{stepbox}

\field{۱. هدف و معرفی مسئله}{
اینجا توضیح دهید مسئله چیست و چرا مهم است.
}

\field{۲. روش/الگوریتم}{
شرح ریاضی یا الگوریتمی روش. برای معادلات:
\[
f(x) = \dots
\]
}

\field{۳. نتایج}{
\begin{center}
\includegraphics[width=.6\textwidth]{images/PLACEHOLDER.png}
\captionof{figure}{توضیح تصویر}
\label{fig:placeholder}
\end{center}
}

\field{۴. جمع‌بندی}{
جمع‌بندی کوتاه یافته‌ها.
}

\end{stepbox}

% در صورت نیاز، جعبه‌های زیر را برای نکات جانبی اضافه کنید:
% \begin{notebox} ... \end{notebox}
% \begin{tipbox} ... \end{tipbox}
% \begin{warnbox} ... \end{warnbox}
% \begin{codebox}[نام فایل] \begin{lstlisting}[language=Python] ... \end{lstlisting} \end{codebox}
